\chapter{Conception du modèle de criticité}

\section{Principes de conception}

La conception proposée repose sur un principe simple : une alerte doit représenter un risque opérationnel clair,
compréhensible et actionnable.

Pour atteindre cet objectif, les règles de conception suivantes sont appliquées :
\begin{itemize}
	\item privilégier les signaux liés aux symptômes de panne plutôt qu'aux seules causes techniques,
	\item expliciter le niveau de criticité de chaque alerte,
	\item associer chaque alerte à un périmètre de responsabilité identifié (service, équipe),
	\item standardiser la structure des notifications pour accélérer le diagnostic initial.
\end{itemize}

\section{Niveaux de criticité et stratégie de notification}

Le modèle visé distingue plusieurs niveaux de criticité, permettant d'adapter la réponse attendue et le canal de
notification.

Sans figer ici une nomenclature définitive, la logique générale est la suivante :
\begin{itemize}
	\item criticité haute : impact utilisateur avéré ou imminent, mobilisation immédiate de l'astreinte,
	\item criticité intermédiaire : dégradation significative à surveiller avec intervention planifiée,
	\item criticité faible : signal d'observabilité utile mais non bloquant à court terme.
\end{itemize}

Cette hiérarchisation doit réduire le bruit et améliorer la qualité des décisions en phase d'incident.

\begin{table}[h]
	\centering
	\caption{Matrice de criticité proposée (version initiale)}
	\label{tab:criticite-initiale}
	\begin{tabular}{|p{0.2\textwidth}|p{0.35\textwidth}|p{0.35\textwidth}|}
		\hline
		\textbf{Niveau} & \textbf{Critère principal} & \textbf{Réponse attendue} \\
		\hline
		Haute & Impact utilisateur avéré ou imminent & Notification immédiate de l'astreinte et traitement prioritaire \\
		\hline
		Intermédiaire & Dégradation significative sans arrêt complet & Investigation planifiée dans une fenêtre courte \\
		\hline
		Faible & Anomalie non bloquante à court terme & Suivi en observation et traitement différé \\
		\hline
	\end{tabular}
\end{table}

\section{Traçabilité, maintenabilité et gouvernance}

Au-delà des seuils d'alarme, la solution doit rester maintenable dans le temps. Le travail de conception intègre
donc des exigences transverses :
\begin{itemize}
	\item traçabilité des règles d'alerting,
	\item conventions homogènes de nommage,
	\item gestion claire du cycle de vie des alarmes,
	\item possibilité d'évolution progressive sans rupture de service.
\end{itemize}

La gouvernance de ces règles est essentielle pour éviter le retour à un système trop bruyant ou incohérent.

\section{Recommandations de déploiement progressif}

Pour sécuriser l'adoption du modèle de criticité, il est recommandé de procéder par itérations :
\begin{itemize}
	\item commencer par un périmètre limité de services à fort enjeu,
	\item mesurer l'effet des changements sur le volume et la qualité des alertes,
	\item ajuster les règles avant extension à l'ensemble de la plateforme.
\end{itemize}

Cette stratégie permet de maîtriser les risques et d'obtenir des retours opérationnels rapides.